% TOP_PROBLEM: {"id": 2836, "title": "Kirby Problem 3.38", "category_id": 7, "category": {"id": 7, "name": "topology", "display_name": "Topology", "description": "Properties preserved under continuous deformations.", "slug": "topology", "order_index": 7, "created_at": "2026-07-31T15:26:25.671Z"}, "status": "open", "problem_number": "KP-3.38", "set_id": 11}
% FIRST_POSTED: 2026-10-01
% ENTRY_KIND: statement_only
% UnsolvedMath ID 2836; source catalogue code KP-3.38
% !TeX program = lualatex
% Definition and sources only; this file is not an LLM solution attempt.
\documentclass[11pt,a4paper]{article}
\usepackage[margin=25mm]{geometry}
\usepackage{fontspec}
\setmainfont{lmroman10-regular.otf}[BoldFont=lmroman10-bold.otf,
 ItalicFont=lmroman10-italic.otf,BoldItalicFont=lmroman10-bolditalic.otf]
\usepackage{amsmath,amssymb,mathtools}
\usepackage{unicode-math}
\setmathfont{latinmodern-math.otf}
\AtBeginDocument{\let\setminus\smallsetminus}
\usepackage{xurl,hyperref}
\hypersetup{colorlinks=true,urlcolor=blue}
\setcounter{secnumdepth}{0}
\setlength{\parindent}{0pt}
\setlength{\parskip}{5pt}
\setlength{\emergencystretch}{3em}
\begin{document}
\section*{Kirby Problem 3.38}\label{2836}
{\small UnsolvedMath ID 2836 · KP-3.38 · Topology · Kirby's Problems in Low-Dimensional Topology}
\subsection{Definitions and mathematical statement}
Use the finitely presented convention for Poincaré duality groups in the
source. A group $G$ is a $\mathrm{PD}_3$ group over $\mathbb Z$ if it is
finitely presented, the trivial module $\mathbb Z$ admits a finite-length
resolution by finitely generated projective $\mathbb ZG$-modules,
its cohomological dimension over $\mathbb Z$ is three, and
\[
 H^i(G;\mathbb ZG)=0\ (i\ne3),\qquad
 H^3(G;\mathbb ZG)\cong\mathbb Z
 \quad\text{as abelian groups}.
\]
The right $G$-action on the last group may be nontrivial; it records the
orientation character $G\to\{\pm1\}$. These conditions express integral
Poincaré duality with the corresponding orientation module, rather than
requiring orientability.

Is every such group isomorphic to the fundamental group of a closed
connected aspherical three-manifold? Aspherical means that the universal
cover is contractible, or equivalently all higher homotopy groups vanish.
Both orientable and nonorientable manifolds are allowed as dictated by the
orientation character.

A realization by a three-dimensional complex satisfying duality alone is
not the conclusion: the complex must be realizable by a genuine closed
manifold. Finite presentation is stated explicitly so that this formulation
is not confused with broader finiteness conventions for duality groups.
\subsection{Short English statement}
Must every finitely presented group satisfying integral three-dimensional Poincaré duality come from a closed aspherical three-manifold?
\subsection{Sources}
\begin{itemize}
\item[S1] ULAM.AI and the UnsolvedMath Contributors, supplied problems.json, record 2836 (KP-3.38), Kirby's Problems in Low-Dimensional Topology. Catalogue identity and source transcription; CC BY 4.0.
\url{https://huggingface.co/datasets/ulamai/UnsolvedMath}.
\item[S2] K3: A New Problem List in Low-Dimensional Topology, Problem 3.38. Expanded formulation of the supplied catalogue statement.
\url{https://math.berkeley.edu/sites/default/files/surv-295-ruberman-watermarked-author-pdf.pdf}.
\end{itemize}
\end{document}
